\cdpchapter{Resumen}

Hacer la compra cada semana plantea una decisi\'on que casi nadie tiene tiempo de resolver bien:
en qu\'e tiendas comprar para gastar menos sin recorrer media ciudad. Los precios var\'ian
notablemente de un supermercado a otro, pero compararlos a mano resulta tedioso y, cuando se hace,
rara vez se tiene en cuenta lo que cuesta desplazarse hasta el establecimiento m\'as barato.

BarGAIN es una aplicaci\'on m\'ovil y web pensada para resolver precisamente eso. A partir de la
lista de la compra del usuario y de su ubicaci\'on, calcula la combinaci\'on de tiendas m\'as
conveniente atendiendo a la vez a tres factores: cu\'anto cuesta la cesta, cu\'anto hay que
desplazarse y cu\'anto tiempo lleva. Cada persona decide a qu\'e da m\'as importancia ---ahorrar,
ir cerca o tardar poco--- y la aplicaci\'on ajusta su recomendaci\'on en consecuencia.

Adem\'as de comparar precios entre supermercados, la aplicaci\'on incorpora tres capacidades que
la distinguen de los comparadores habituales: propone la ruta de compra m\'as ventajosa repartida
entre varias tiendas; permite crear la lista a partir de una foto de un ticket o de una nota
escrita a mano, sin teclear; y ofrece un asistente conversacional al que se puede preguntar en
lenguaje corriente por precios, ofertas o por la mejor manera de organizar la compra. Un portal
web complementario da cabida al peque\~no comercio local, que puede publicar sus precios y
promociones para competir en igualdad de condiciones con las grandes cadenas.

El proyecto se ha llevado a cabo recorriendo el ciclo completo de la ingenier\'ia del software
---an\'alisis de requisitos, dise\~no, implementaci\'on, pruebas y despliegue---. El resultado es
un prototipo plenamente funcional, accesible de forma p\'ublica y verificado mediante pruebas
automatizadas y ensayos en un dispositivo real. La memoria recoge tambi\'en las decisiones de
dise\~no adoptadas, las limitaciones de esta primera versi\'on y las l\'ineas de mejora previstas
para convertir BarGAIN en un producto a mayor escala.

\cdpchapter{Abstract}

Doing the weekly grocery shopping involves a decision few people have time to get right: which
shops to visit in order to spend less without driving halfway across town. Prices vary widely from
one supermarket to another, but comparing them by hand is tedious and, even when done, it rarely
accounts for the cost of travelling to the cheapest store.

BarGAIN is a mobile and web application designed to solve exactly that. Starting from the user's
shopping list and location, it works out the most convenient combination of shops by weighing
three factors at once: how much the basket costs, how far one has to travel, and how long it
takes. Each person chooses what matters most ---saving money, staying nearby, or spending little
time--- and the application adapts its recommendation accordingly.

Beyond comparing supermarket prices, the application offers three capabilities that set it apart
from conventional comparators: it suggests the most advantageous shopping route across several
shops; it lets users build their list from a photo of a receipt or a handwritten note, with no
typing required; and it provides a conversational assistant that answers everyday questions about
prices, offers, or how best to organise the shopping. A complementary web portal welcomes small
local businesses, which can publish their prices and promotions to compete on equal terms with the
large chains.

The project was carried out following the full software engineering lifecycle ---requirements
analysis, design, implementation, testing, and deployment---. The outcome is a fully functional,
publicly accessible prototype, validated through both automated tests and trials on a real device.
The report also documents the design decisions made, the limitations of this first version, and
the planned improvements toward turning BarGAIN into a larger-scale product.
